\answer{OA33}
初始化历史峰值 $H=p_1$、峰值下标 $i_H=1$、答案 $D=0$。对 $j=2,\ldots,n$ 依次执行：先计算 $d=1-p_j/H$；若 $d>D$，更新 $D=d$ 并保存区间 $(i_H,j)$；然后若 $p_j>H$，更新 $H=p_j$、$i_H=j$。若始终没有正跌幅，就返回 $D=0$，且无需返回正回撤区间。

处理 $j$ 之前，$H=\max_{i<j}p_i$。对固定的正数 $p_j$，函数 $1-p_j/p_i$ 随正数 $p_i$ 增大而增大，所以以 $j$ 为终点的最大跌幅恰为 $1-p_j/H$。算法对每个终点都检查了它的最优起点，再与零取最大值，因此得到全局最大回撤。先计算再更新峰值保证了起点严格早于终点。

每个价格只处理一次，时间为 $O(n)$；除了几个数值和下标外不需要额外存储，辅助空间为 $O(1)$。使用严格的 $d>D$ 更新会保留最早出现的最大回撤终点，峰值相等时不更新则保留最早峰值下标。

回撤的分母必须是此前峰值。例如价格从 $100$ 跌到 $60$，回撤为 $40\%$。表达式 $H/p_j-1$ 算的是从低点恢复到峰值所需的涨幅，在本例中为 $2/3$，不是回撤。若用累计峰值 $H_j=\max_{i\le j}p_i$ 表示，同一答案可写为
\[
 \boxed{D=\max_{1\le j\le n}\left(1-\frac{p_j}{H_j}\right)}.
\]
